% 三个观测对中分别选择一个；所示选择向量为 (1,0,1)。
\begin{tikzpicture}[x=1mm,y=1mm,font=\kaishu\small,
 obs/.style={book node,minimum width=18mm,minimum height=8mm,align=center},
 chosen/.style={obs,fill=FigureInput!18,draw=FigureInput},
 spare/.style={obs,fill=BookPaper},
 information output/.style={book node,minimum width=23mm,minimum height=12mm,align=center,fill=FigureModel!12}]
 \node[anchor=east] at (-14,6) {备选0};
 \node[anchor=east] at (-14,-6) {备选1};
 \node[spare] (a0) at (0,6) {$z_{1,0}$};
 \node[chosen] (a1) at (0,-6) {$z_{1,1}$};
 \node[chosen] (b0) at (32,6) {$z_{2,0}$};
 \node[spare] (b1) at (32,-6) {$z_{2,1}$};
 \node[spare] (c0) at (64,6) {$z_{3,0}$};
 \node[chosen] (c1) at (64,-6) {$z_{3,1}$};
 \node at (0,18) {$B_1=1$};
 \node at (32,18) {$B_2=0$};
 \node at (64,18) {$B_3=1$};
 \node[chosen] (sa) at (0,-29) {$z_{1,1}$};
 \node[chosen] (sb) at (32,-29) {$z_{2,0}$};
 \node[chosen] (sc) at (64,-29) {$z_{3,1}$};
 \node[anchor=east] at (-14,-29) {$S_{\bm B}$};
 \draw[book arrow] (a1)--(sa);
 \draw[book arrow] (b0.east)--(45,6)--(45,-18)--(32,-18)--(sb.north);
 \draw[book arrow] (c1)--(sc);
 \node[information output] (alg) at (105,-29) {学习算法\\$W=A(S_{\bm B},U)$};
 \draw[book arrow] (sc.east)--(alg.west);
 \node[align=center,font=\kaishu\footnotesize] at (106,1) {给定全部备选观测，\\衡量输出对$\bm B$的信息量};
\end{tikzpicture}
